% FTAD report — preamble injected via pandoc's -H option.
%
% The report specification is Tahoma 12 pt at 1.5 line spacing on A4, with a
% "Page X of Y" footer. That rules out the IEEEtran two-column layout, so the
% document is a single-column `report` class with IEEE-style numbered citations
% written inline in the sources. IEEE refers to the citation convention here,
% not to the IEEEtran class.

\usepackage{fontspec}

% Tahoma is resolved through fontconfig. On this host that is wine-tahoma,
% which is metric-compatible with the Microsoft original.
\setmainfont{Tahoma}
\setsansfont{Tahoma}
\setmonofont{Courier New}

% 1.5 line spacing, applied document-wide.
\linespread{1.5}

\usepackage{fancyhdr}
\usepackage{lastpage}

\pagestyle{fancy}
\fancyhf{}
\fancyfoot[C]{Page \thepage\ of \pageref{LastPage}}
\renewcommand{\headrulewidth}{0pt}
\renewcommand{\footrulewidth}{0pt}

% Headings carry their own numbers in the source text ("Chapter 4 — ...",
% "4.1 Introduction"), so LaTeX must not add a second set. secnumdepth -1 stops
% the table of contents from prepending its own index: at 0 the chapters were
% still numbered, so the TOC read "6 Chapter 4 — Testing Methodology" because
% Abstract and Introduction occupied slots 1 and 2. Sections keep their own
% numbers in the heading text and remain listed.
\setcounter{secnumdepth}{-1}

\setcounter{tocdepth}{2}

% Tables at 10 pt on single spacing. Body prose keeps 12 pt / 1.5 as specified;
% tables and listings do not, because a table row set at 1.5 spacing is 1.5x
% taller than it needs to be and a wide table then wraps into unreadable slabs.
\usepackage{etoolbox}
\usepackage{setspace}
\AtBeginEnvironment{longtable}{\scriptsize\singlespacing}
\AtBeginEnvironment{tabular}{\scriptsize\singlespacing}

% Keep code listings from running off the page, and set them single-spaced for
% the same reason as the tables.
\usepackage{fvextra}
\DefineVerbatimEnvironment{Highlighting}{Verbatim}{breaklines,commandchars=\\\{\}}
\AtBeginEnvironment{Highlighting}{\linespread{1}\selectfont\footnotesize}
\AtBeginEnvironment{verbatim}{\linespread{1}\selectfont\footnotesize}

% pandoc's template loads parskip, whose default \parskip is 0.5\baselineskip --
% about 11 pt at 12 pt with 1.5 spacing. On a page of short paragraphs that is
% 70-80 mm of whitespace, and it was the main reason the first draft ran to 99
% pages. Paragraphs stay separated, but compactly.
\setlength{\parskip}{2pt plus 1pt minus 1pt}
\setlength{\parindent}{0pt}

% Tighten heading spacing for the same reason. Chapters still start on a new
% page, which is conventional for a report.
\usepackage{titlesec}
\titlespacing*{\chapter}{0pt}{6pt}{12pt}
\titlespacing*{\section}{0pt}{8pt}{3pt}
\titlespacing*{\subsection}{0pt}{6pt}{2pt}

% Chapters flow continuously rather than each forcing a page break. The report
% body is capped at 40 pages, and nine chapter boundaries ending mid-page were
% costing roughly four pages of blank paper. Rather than rebuild the chapter
% machinery with titlesec's \titleclass (which looped), the page break is
% patched out of report.cls directly, and the plain page style is patched away
% so the "Page X of Y" footer survives on chapter-opening pages. Back matter
% still starts on a new page, injected as raw LaTeX by assemble.py.
\makeatletter
\patchcmd{\chapter}{\if@openright\cleardoublepage\else\clearpage\fi}{}{}{}
\patchcmd{\chapter}{\thispagestyle{plain}}{}{}{}
\makeatother


% Long inline option strings ("-fsanitize=address,undefined
% -fno-omit-frame-pointer") cannot break at a space inside \texttt and were
% running up to 310 pt off the page. [htt] lets them hyphenate, and the
% emergency stretch absorbs the remainder without visible looseness.
\usepackage[htt]{hyphenat}
\setlength{\emergencystretch}{4em}

% pandoc loads xurl, which lets a URL break at any character, so the reference
% list read "https://cppcheck.sou / rceforge.io/manual.pdf". Restricting the
% break points to punctuation keeps URLs legible; every entry still fits, since
% the longest unbroken segment is 22 characters. \AtBeginDocument is used
% because xurl redefines \UrlBreaks when it loads, after this file.
% longtable centres tables narrower than the text block, so the short two- and
% three-column tables floated in the middle of the page while the body text is
% left-aligned. Aligning them to the left margin is the convention in the rest
% of the report.
\AtBeginDocument{%
  \setlength{\LTleft}{0pt}%
  \setlength{\LTright}{\fill}%
}

\AtBeginDocument{\def\UrlBreaks{\do\/\do\-\do\.\do\_\do\:}}

\usepackage{microtype}
